\scanpage{46}
% archive sheet number: 45
$A$ variété abélienne en car $p$. On a, pour la partie de $\pi_1(A)$ première à $p$
\[
\pi_1^{\ne p}(A)\simeq\varprojlim_{\substack{n\\\text{\unsure{arbitr.|arbitre|entier}}\\(n,p)=1}}
\deleted{\unsure{a|n}}A^{(n)}
\]
où $\deleted{\unsure{n|A}}A^{(n)}$ est le noyau de $A\xrightarrow{n}A$, et où $A^{(mn)}\unsure{-|\to} A^{(n)}$ est la multiplication par $m$.
\[
\boxed{\pi_1^\ell(A)\simeq\varprojlim_{\deleted{\unsure{\ell|n}}k}A^{(\ell^k)}
\simeq T_\ell(A)=\operatorname{Hom}(\underline{\mathbb Q}_\ell/\mathbb Z_\ell,A)}
\]
module de Tate

$\ell$ étant un nb premier $\ne p$. Ce résultat provient du fait qu'un \uline{revêtement sans \unsure{ramifi.|ramif.} d'une variété abélienne est une variété abélienne}

D'après Kummer, on a d'autre part
\[
\deleted{\unsure{\mathrm{Dual}|\mathrm{Du}}}\ \operatorname{Hom}(\pi_1^\ell(A),\unsure{U|\mu}_\ell)\simeq
\text{sous-groupe de }\unsure{\mathscr D|\mathscr P}(A)=H^1(A,\underline O_A^{\unsure{\times|*}})
\]
des éléments d'ordre une puissance de $\ell$.

Si on admet que \uline{$A$ n'a pas de torsion}, on obtient donc
\[
\operatorname{Hom}(\pi_1^\ell(A),U_\ell)\simeq A^{*(\ell^\infty)},\qquad\text{i.e.}
\]
\[
\boxed{T_\ell(A)\times T_\ell(A^*)\longrightarrow T_\ell(U_\ell)}
\]
accouplement à valeurs dans le $\mathbb Z_\ell$-module libre de rang $1$ $T_\ell\unsure{(U_\ell)|U_\ell)}$, \unsure{et|cet} accouplement qui est une dualité.
\[
\begin{aligned}
&T_\ell(A)/\ell^kT_\ell(A)\times T_\ell(A^*)/\ell^kT_\ell\unsure{(A^*)|A^*)}
\to T_\ell(U_\ell)/\unsure{(\ell^kT_\ell U_\ell)|\ell^kT_\ell(U_\ell)}\\
&A^{(\ell^k]}\times A^{*(\ell^k)}\longrightarrow U^{(\ell^k)}
\end{aligned}
\]
